% Teil II: durchgerechneter Weg, Teile a–e.

\subsection*{Gegeben und gesucht}
Gegeben: $r_1=\qty{10}{cm}$, $z_0=\qty{30}{cm}$, $r_2=\qty{2.5}{mm}$,
$l=\qty{10}{cm}$, $\eta=\qty{100}{mPa\,s}=\qty{0.1}{Pa\,s}$,
$\rho=\qty{920}{kg/m^3}$, $g=\qty{9.81}{m/s^2}$.
Gesucht: $v_2(z)$ (a), DGL+AB für $z(t)$ (b), Halbentleerungszeit ideal (c),
Anfangs-$\bar u_2$ (d), Halbentleerungszeit mit Reibung (e).

\subsection*{Ansatz}
(a) Bernoulli/Torricelli aufs offene Gefäß angewandt; wegen $r_2\ll r_1$ ist
die Oberflächengeschwindigkeit vernachlässigbar. (b)–(c) Kontinuität
$A_1\dot z=-A_2 v_2$ mit $v_2=\sqrt{2gz}$; Trennung der Variablen. (d)–(e)
Hagen-Poiseuille statt Torricelli; mittlere Geschwindigkeit $\bar u\propto z$,
daraus lineare DGL und Exponentialgesetz.

\subsection*{Rechnung}
\textbf{a) Ausflussgeschwindigkeit.} Torricelli mit aktueller Höhe $z$:
\begin{equation}
  v_2=\sqrt{2gz}\;.
\end{equation}

\textbf{b) DGL für $z(t)$.} Einsetzen in die Kontinuität:
\begin{equation}
  -\pi r_1^2\,\dot z \;=\; \pi r_2^2\sqrt{2gz}
  \;\Rightarrow\;
  \boxed{\;\dot z=-\frac{r_2^2}{r_1^2}\sqrt{2gz}\;},\qquad z(0)=z_0=\qty{0.30}{m}\;.
\end{equation}

\textbf{c) Halbentleerungszeit (ideal).} Trennung der Variablen:
\begin{equation}
  \frac{\mathrm dz}{\sqrt z}=\frac{\mathrm{d}z}{\sqrt z}=\sqrt{\frac{r_2^2}{r_1^2}\cdot 2g}\,(-\mathrm dt)
  \;\Rightarrow\;
  2\sqrt z=2\sqrt{z_0}-\frac{r_2^2}{r_1^2}\sqrt{2g}\,t\;.
\end{equation}
Für $z=z_0/2$:
\begin{equation}
  t_{\tfrac12}
  =\frac{2\left(\sqrt{z_0}-\sqrt{z_0/2}\right)}{\frac{r_2^2}{r_1^2}\sqrt{2g}}\;.
\end{equation}
Zahlen: $\frac{r_2^2}{r_1^2}=\left(\frac{0.0025}{0.10}\right)^2=6.25\times10^{-4}$,
$\sqrt{z_0}=0.5477$, $\sqrt{z_0/2}=0.3873$, $\sqrt{2g}=4.429$:
\begin{equation*}
  t_{\tfrac12}=\frac{2(0.5477-0.3873)}{6.25\times10^{-4}\cdot4.429}
  \approx \qty{116}{s}\approx\qty{1.9}{min}.
\end{equation*}

\textbf{d) Mittlere Ausflussgeschwindigkeit bei $t=0$ (mit Reibung).}
Hagen-Poiseuille, $\Delta p=\rho g z_0$ (hydrostatischer Druck der Säule):
\begin{equation}
  \bar u(t{=}0)=\frac{r_2^2}{8\eta l}\,\rho g z_0
  =\frac{(0.0025)^2\cdot 920\cdot 9.81\cdot 0.30}
        {8\cdot 0.1\cdot 0.10}
  \approx \qty{0.21}{m/s}\;.
\end{equation}

\textbf{e) Warum quadratische Terme entfallen; lineare DGL und Halbzeit.}
Dynamischer Druck $\tfrac\rho2 \bar u^2\approx \tfrac{920}{2}(0.21)^2
\approx\qty{20}{Pa}$ gegen hydrostatisch $\rho g z_0\approx\qty{2710}{Pa}$ —
Faktor $\sim100$ kleiner, also $\tfrac\rho2 v^2$ vernachlässigbar. Dann gilt
linear $\bar u(z)=\frac{r_2^2}{8\eta l}\rho g\, z\equiv C z$ mit
$C=\frac{\bar u(t{=}0)}{z_0}=\frac{0.21}{0.30}=\qty{0.71}{s^{-1}}$.
Kontinuität ergibt
\begin{equation}
  \dot z=-\frac{r_2^2}{r_1^2}\,C\,z
  \;\Rightarrow\;
  z(t)=z_0\,e^{-k t},\qquad
  k=\frac{r_2^2}{r_1^2}\,C=6.25\times10^{-4}\cdot 0.71
    \approx 4.4\times10^{-4}\,\mathrm{s^{-1}}\;.
\end{equation}
Halbentleerung:
\begin{equation*}
  t_{\tfrac12}=\frac{\ln 2}{k}\approx\frac{0.693}{4.4\times10^{-4}}
  \approx \qty{1570}{s}\approx\qty{26}{min}.
\end{equation*}

\begin{ergebnis}
\begin{equation*}
  v_2=\sqrt{2gz}\;;\quad
  \dot z=-\frac{r_2^2}{r_1^2}\sqrt{2gz},\ z(0)=z_0\;;\quad
  t_{\tfrac12}^{\mathrm{ideal}}\approx\qty{116}{s}\;;\quad
  \bar u(0)\approx\qty{0.21}{m/s}\;;\quad
  t_{\tfrac12}^{\mathrm{Reib}}\approx\qty{26}{min}.
\end{equation*}
\end{ergebnis}

\begin{kontrolle}
\textbf{Einheiten:} $t_{\tfrac12}$
$\to\sqrt{m}/(m^2/m^2\cdot \sqrt{m/s^2}\,\mathrm{s})$, richtig $\mathrm{s}$;
$k$ in $\mathrm{s^{-1}}$, $\ln2/k$ in $\mathrm{s}$. \checkmark

\textbf{Grenzfall/Querprobe:} Ideal $t_{\tfrac12}\propto\sqrt{z_0}$ (längere
Säule → länger), reibend $t_{\tfrac12}\propto 1/\eta$ (zäher → viel länger);
mit Reibung deutlich länger als ideal — physikalisch sinnvoll.

\textbf{Größenordnung:} Ohne Reibung ~2 min, mit dem dünnen Rohr ~26 min —
plausibel für Ausströmen durch eine 2.5-mm-Kapillare.
\end{kontrolle}

\begin{merksatz}
Eine Entleerungsaufgabe läuft immer auf eine DGL aus Kontinuität
$A_1\dot z=-A_2 v_2(z)$ hinaus; Torricelli liefert eine $\sqrt z$-DGL
(Integral über $1/\sqrt z$), Hagen-Poiseuille eine lineare (Exponentiallösung).
Ob der dynamische Term zählt, entscheidet der Vergleich
$\rho v^2/2$ gegen $\rho g z$.
\end{merksatz}
